\section{Conclusion}
\label{sec:conclusion}

We presented a sheaf-theoretic framework for interpretable knowledge graph
reasoning. At a single scale (the type graph), cellular sheaf cohomology
gives a quantitative inconsistency signal: $\dim H^1$ rises by $+53$ when
76 conflicts are injected into a clean graph, with roughly 70\% of
injected conflicts producing linearly independent cohomology classes
(Theorem~\ref{thm:conflict_sensitivity_full}). On top of this, ontology-typed
attention drives invalid attention weight to zero by construction, and
proof-guided generation emits a derivation alongside each prediction. Link
prediction performance on FB15K-237 is on par with KG-BERT, KGT5, and
SimKGC; our contribution is the structural and verifiable layer that none
of these baselines provide.

\paragraph{From a controlled environment to a deployable system.}
The results above are established in a deliberately controlled
environment: hand-built Ologs of $\leq 100$ types, a 7-type e-commerce
schema, synthetic label priors, and single-seed runs. This is the right
setting for isolating the mechanisms---it lets us measure $H^1$, the
(B)/(D) separation, and the topology contrast without confounds---but it
is not yet a system competitive with state-of-the-art neural KG
reasoners. Four research blockers stand between the present work and
that goal; a fifth direction points outward, generalizing the framework
beyond knowledge graphs altogether.
\begin{enumerate}
\item \textbf{Cohomology at entity scale.} Cellular sheaf cohomology with
non-trivial restriction maps runs in $O((nd)^3)$ (dense rank on the
boundary operators). This is why $H^1$ is reported on the type graph
($n \leq 100$) and only the cheaper $O(n+m)$ Betti number on the
entity graph. \emph{Blocker:} a sparse-Krylov sheaf-cohomology
algorithm exploiting block structure. Until it exists, the calibrated
$H^1$ signal cannot be computed directly on a full KG.
\item \textbf{From synthetic priors to trained models.} The (B$'$)/(D)
trajectory results use hand-crafted label priors and a structural audit
of the production attention layer, not an end-to-end trained language
model with a learned label head. \emph{Blocker:} a trained projection
head from attention output to label distribution, which would also
unlock the four-cell (A / B / D / B+D) ablation needed to settle what
(B) contributes when (D) is present (Appendix~\ref{app:b_vs_d}).
\item \textbf{From hand-built Ologs to learned ontological structure.}
Every Olog in this paper is hand-authored, and the entity embeddings and
tokenization are kept deliberately simple. A deployable system needs
ontologies induced from text or data at scale, \emph{and} embeddings and
a tokenization that faithfully encode the Olog's type structure rather
than merely fit a link-prediction objective. \emph{Blocker:} both are
open problems. Reliable automated Olog extraction would remove the
hand-authoring bottleneck; ontology-faithful representation
learning---in our view the harder of the two, and a central open problem
for this line of work---would let the type structure \emph{shape} the
embedding geometry rather than be imposed only as an external attention
mask, as it is here.
\item \textbf{Multi-scale unification.} We report two disconnected
topological measurements---sheaf $H^1$ on the type graph and graph
$b_1$ on the entity graph. \emph{Blocker:} a sheaf morphism across
scale, or a persistent sheaf cohomology over a filtration that merges
relations within an $\epsilon$-radius, unifying both into a single
multi-scale statement that subsumes the present work as the
single-scale special case.
\item \textbf{Beyond knowledge graphs.} The framework takes any
Olog-typed graph as input, which invites instantiation well beyond KG
reasoning---though the targets are not equally close. A database
entity-relationship diagram is a near-direct fit: an ERD \emph{is} a
schema, its instances \emph{are} typed graphs, and referential-integrity
or constraint violations are the kind of inconsistency the sheaf $H^1$
already measures. Causal diagrams and Bayesian
networks~\cite{pearl2009causality} are a genuine stretch: their edges
carry interventional and conditional-dependence semantics, and the
inconsistency that matters there---d-separation violations, incoherent
probability mass---is not the same object as restriction maps
disagreeing. \emph{Blocker:} a sheaf whose stalks and restriction maps
encode distributional, not merely type, structure---with a corresponding
cohomological notion of inconsistency---has to be worked out before the
method transfers to the probabilistic targets.
\end{enumerate}

\paragraph{Structure before correlation.} This pattern is not new.
Pearl's program for causal inference rests on one premise: intervention
and counterfactual reasoning cannot be recovered from a joint
distribution alone---they require a structural model supplied from
outside the data~\cite{pearl2009causality}. Our claim is the same move
along a different axis: consistency and verifiability cannot be
recovered from a link-prediction objective alone; they require the
ontology, supplied as an explicit object the system reasons over.
Accuracy measures how well a model fits the data it was given and is
structurally silent on whether a prediction respects the schema behind
that data. Making the Olog explicit is what gives that question a place
to live.

\paragraph{Toward auditable reasoning.}
The objective behind this work is not a higher MRR. Knowledge graphs
increasingly feed decisions in medicine, finance, and science, where a
confident but ontologically incoherent prediction is not a ranking error
but a safety failure---and one that accuracy metrics are structurally
blind to. Our aim is to give KG reasoning a layer that can be
\emph{audited}: a quantitative inconsistency signal that says where in
the graph coherence breaks down, and a generation pipeline in which every
prediction either carries a derivation or is flagged as unfounded. That
shifts the question a practitioner can ask of a model from ``is it
accurate on average?'' to ``can this specific prediction be trusted, and
can I see why?''. We see the sheaf-cohomological and two-locus
constructions here as a first, deliberately small step toward reasoning
systems that are verifiable by construction rather than trusted by
reputation.
